% Compile with: pdflatex lecture_28.tex
% Requires: amsmath, amssymb, amsthm, mathtools, xcolor, mdframed,
%           enumitem, hyperref, pagecolor, microtype

\documentclass[11pt]{article}
\usepackage[margin=1in]{geometry}
\usepackage{amsmath, amssymb, amsthm}
\usepackage{mathtools}
\usepackage{xcolor}
\usepackage{mdframed}
\usepackage{enumitem}
\usepackage{hyperref}
\usepackage{pagecolor}
\usepackage{microtype}

% ---------- Colour palette ----------
\definecolor{cream}{HTML}{FFFDF5}
\definecolor{defblue}{HTML}{1A4A7A}
\definecolor{thmgreen}{HTML}{1A5C3A}
\definecolor{warnAmber}{HTML}{7A4A00}
\definecolor{dangerRed}{HTML}{7A1A1A}
\definecolor{boxbluebg}{HTML}{EBF3FB}
\definecolor{boxgreenbg}{HTML}{EBF7EF}
\definecolor{boxamberbg}{HTML}{FDF5E6}
\definecolor{boxredbg}{HTML}{FBEBEB}
\pagecolor{cream}

% ---------- Theorem styles ----------
\newtheoremstyle{defstyle}    {8pt}{8pt}{\normalfont}{0pt}{\bfseries\color{defblue}} {.}{0.5em}{}
\newtheoremstyle{thmstyle}    {8pt}{8pt}{\normalfont}{0pt}{\bfseries\color{thmgreen}}{.}{0.5em}{}
\newtheoremstyle{notstyle}    {8pt}{8pt}{\normalfont}{0pt}{\bfseries\color{warnAmber}}{.}{0.5em}{}
\newtheoremstyle{cautionstyle}{8pt}{8pt}{\normalfont}{0pt}{\bfseries\color{dangerRed}}{.}{0.5em}{}

\theoremstyle{defstyle}
\newtheorem{definition}{Definition}[section]

\theoremstyle{thmstyle}
\newtheorem{theorem}{Theorem}[section]
\newtheorem{corollary}[theorem]{Corollary}
\newtheorem{lemma}[theorem]{Lemma}

\theoremstyle{notstyle}
\newtheorem*{remark}{Remark}
\newtheorem*{example}{Example}
\newtheorem*{addendum}{Addendum}

\theoremstyle{cautionstyle}
\newtheorem*{caution}{Caution}

% ---------- Boxes ----------
\newmdenv[
  backgroundcolor=boxbluebg, linecolor=defblue, linewidth=2pt,
  topline=false, rightline=false, bottomline=false,
  innerleftmargin=10pt, innerrightmargin=10pt,
  innertopmargin=6pt, innerbottommargin=6pt,
  skipabove=8pt, skipbelow=8pt
]{defbox}

\newmdenv[
  backgroundcolor=boxgreenbg, linecolor=thmgreen, linewidth=2pt,
  topline=false, rightline=false, bottomline=false,
  innerleftmargin=10pt, innerrightmargin=10pt,
  innertopmargin=6pt, innerbottommargin=6pt,
  skipabove=8pt, skipbelow=8pt
]{thmbox}

\newmdenv[
  backgroundcolor=boxamberbg, linecolor=warnAmber, linewidth=2pt,
  topline=false, rightline=false, bottomline=false,
  innerleftmargin=10pt, innerrightmargin=10pt,
  innertopmargin=6pt, innerbottommargin=6pt,
  skipabove=8pt, skipbelow=8pt
]{notebox}

\newmdenv[
  backgroundcolor=boxredbg, linecolor=dangerRed, linewidth=2pt,
  topline=false, rightline=false, bottomline=false,
  innerleftmargin=10pt, innerrightmargin=10pt,
  innertopmargin=6pt, innerbottommargin=6pt,
  skipabove=8pt, skipbelow=8pt
]{cautionbox}

% ---------- Macros ----------
\newcommand{\Z}{\mathbb{Z}}
\newcommand{\N}{\mathbb{N}}
\newcommand{\R}{\mathbb{R}}
\newcommand{\card}[1]{\lvert #1 \rvert}

\begin{document}

\begin{center}
  {\Large\bfseries Lecture 28 --- Cardinality}\\[4pt]
  {\normalsize CS 173: Discrete Structures \quad$\cdot$\quad Comparing and Counting Sizes of Sets}\\[2pt]
  {\small\color{gray} University of Illinois at Urbana-Champaign}
\end{center}
\vspace{1em}\hrule\vspace{1.5em}

% ============================================================
\section{Review and the Empty-Set Addendum}
% ============================================================

Cardinality comparisons must work for infinite sets too, so we phrase them in
terms of functions, not counts.

\begin{defbox}
\begin{definition}[Cardinality comparison]
For any sets $A$ and $B$:
\begin{enumerate}[label=\textnormal{(\roman*)}, leftmargin=2em, itemsep=2pt]
  \item $\card{A} \leq \card{B} \iff \exists f \;(f : A \to B \,\land\, f \text{ injective})$.
  \item If $B \neq \emptyset$:
        $\card{A} \geq \card{B} \iff \exists f \;(f : A \to B \,\land\, f \text{ surjective})$.
\end{enumerate}
\end{definition}
\end{defbox}

\begin{notebox}
\begin{addendum}[Empty target]
We stipulate $\card{A} \geq \card{\emptyset}$ for all $A$.
\end{addendum}
\end{notebox}

\paragraph{Why the side condition?}
If $B = \emptyset$ and $A \neq \emptyset$, no function $A \to \emptyset$
exists, so the surjection definition would falsely give
$\card{A} \not\geq \card{\emptyset}$. We patch this by fiat. The injection
definition needs no such fix: the empty function $\emptyset \to B$ is
vacuously injective, so $\card{\emptyset} \leq \card{B}$ comes for free.

% ============================================================
\section{Order-Theoretic Properties of $\leq$}
% ============================================================

\begin{thmbox}
\begin{theorem}[Reflexivity]
$\forall A \;(\card{A} = \card{A})$.
\end{theorem}
\end{thmbox}
\begin{proof}
$\mathrm{id}_A : A \to A$ is a bijection.
\end{proof}

\begin{thmbox}
\begin{theorem}[Duality]
$\forall A\, \forall B \;(\card{A} \leq \card{B} \iff \card{B} \geq \card{A})$.
\end{theorem}
\end{thmbox}

\begin{thmbox}
\begin{theorem}[Transitivity]
$\forall A\, \forall B\, \forall C \;\bigl((\card{A} \leq \card{B} \,\land\, \card{B} \leq \card{C}) \Rightarrow \card{A} \leq \card{C}\bigr)$.
\end{theorem}
\end{thmbox}
\begin{proof}
If $f : A \to B$ and $g : B \to C$ are injective, so is $g \circ f$: from
$g(f(a_1)) = g(f(a_2))$, injectivity of $g$ then $f$ yields $a_1 = a_2$.
\end{proof}

\begin{thmbox}
\begin{theorem}[Antisymmetry --- Cantor--Schr\"oder--Bernstein]
$\forall A\, \forall B \;\bigl((\card{A} \leq \card{B} \,\land\, \card{A} \geq \card{B}) \Rightarrow \card{A} = \card{B}\bigr)$.
\end{theorem}
\end{thmbox}

\begin{notebox}
\begin{remark}
For finite sets, antisymmetry is immediate by counting; for infinite sets, a
bijection has to be \emph{built} from the two given functions. This is the
content of Cantor--Schr\"oder--Bernstein, proved separately.
\end{remark}
\end{notebox}

\subsection*{Summary}

\begin{center}
\renewcommand{\arraystretch}{1.4}
\begin{tabular}{|l|l|}
\hline
\textbf{Property} & \textbf{Statement} \\
\hline
Reflexive     & $\card{A} = \card{A}$ \\
Duality       & $\card{A} \leq \card{B} \iff \card{B} \geq \card{A}$ \\
Transitive    & $\card{A} \leq \card{B} \,\land\, \card{B} \leq \card{C} \Rightarrow \card{A} \leq \card{C}$ \\
Antisymmetric & $\card{A} \leq \card{B} \,\land\, \card{A} \geq \card{B} \Rightarrow \card{A} = \card{B}$ \\
\hline
\end{tabular}
\end{center}

Reflexivity, antisymmetry, and transitivity make $\leq$ a partial order on
cardinalities. Duality just says $\geq$ is $\leq$ read backwards.

% ============================================================
\section{Inverses}
% ============================================================

To prove duality, we need: every injection has a surjective left inverse, and
every surjection has an injective right inverse.

\begin{thmbox}
\begin{theorem}[Left and right inverses]
For any $X, Y$ and $f : X \to Y$:
\begin{enumerate}[label=\textnormal{(\Roman*)}, leftmargin=2.5em, itemsep=2pt]
  \item If $f$ is injective, then $\exists g \;(g : Y \to X \,\land\, g \text{ surjective} \,\land\, g \circ f = \mathrm{id}_X)$.
  \item If $f$ is surjective, then $\exists g \;(g : Y \to X \,\land\, g \text{ injective} \,\land\, f \circ g = \mathrm{id}_Y)$.
\end{enumerate}
\end{theorem}
\end{thmbox}

The names ``left'' and ``right'' refer to where $g$ sits in the composition.
Applied to duality: an injection $A \to B$ produces a surjection $B \to A$
via~(I), giving $\card{B} \geq \card{A}$; conversely~(II) turns a surjection
into an injection.

\begin{cautionbox}
\begin{caution}[Empty-domain corner case]
Part~(I) fails as stated when $X = \emptyset$ and $Y \neq \emptyset$: the
empty function $\emptyset \to Y$ is vacuously injective, but no
$g : Y \to \emptyset$ exists. Either assume $X \neq \emptyset$ or handle the
case separately, as with the empty-set addendum in Section~1.
\end{caution}
\end{cautionbox}

\subsection*{The two-sided inverse}

When $f$ is bijective, both parts apply, and the resulting inverse is unique
and itself a bijection.

\begin{thmbox}
\begin{corollary}[Inverse of a bijection]
If $f : X \to Y$ is a bijection, then
\[
  \exists ! g \;\bigl( g : Y \to X
    \,\land\, g \text{ bijection}
    \,\land\, g \circ f = \mathrm{id}_X
    \,\land\, f \circ g = \mathrm{id}_Y \bigr).
\]
\end{corollary}
\end{thmbox}

\begin{defbox}
\begin{definition}[Inverse $f^{-1}$]
The unique $g$ above is called \textbf{the inverse} of $f$, written
$f^{-1} : Y \to X$.
\end{definition}
\end{defbox}

The jump from $\exists$ (one-sided) to $\exists !$ (two-sided) is real: a
non-surjective injection has many left inverses, since values on $Y \setminus
f(X)$ can be chosen freely. A bijection leaves no such freedom.

% ============================================================
\section{Categorical Vocabulary: Monic, Epic, Iso}
% ============================================================

The same three notions can be characterised purely by how $f$ composes with
other functions --- no reference to elements.

\begin{thmbox}
\begin{theorem}[Cancellation characterisations]
For any $X, Y$ and $f : X \to Y$:
\begin{enumerate}[label=\textnormal{(\arabic*)}, leftmargin=2em, itemsep=4pt]
  \item $f$ injective $\iff$ $f$ is \textbf{monic}: \;
        $\forall W,\, \forall h_1, h_2 : W \to X,\;
         f \circ h_1 = f \circ h_2 \Rightarrow h_1 = h_2$.
  \item $f$ surjective $\iff$ $f$ is \textbf{epic}: \;
        $\forall Z,\, \forall k_1, k_2 : Y \to Z,\;
         k_1 \circ f = k_2 \circ f \Rightarrow k_1 = k_2$.
  \item $f$ bijection $\iff$ $f$ is an \textbf{isomorphism}: $f$ has a two-sided inverse.
\end{enumerate}
\end{theorem}
\end{thmbox}

\begin{notebox}
\begin{remark}[Why this matters]
Mono- (one), epi- (onto), iso- (equal) match the elementary names. The point
of the rephrasing: injective/surjective/bijective look \emph{inside} $X$ and
$Y$; monic/epic/iso only mention compositions. This generalises to settings
where objects don't have ``elements'' --- the entry point to category theory.
\end{remark}
\end{notebox}

% ============================================================
\section{Finite Sets}
% ============================================================

So far ``$\card{X}$'' has been a comparison device, not a number. For finite
sets we can finally pin it down to a natural number.

\begin{defbox}
\begin{definition}[Finite set]
$X$ is \textbf{finite} iff $\exists n \in \N \;(\card{X} = \card{n})$.
\end{definition}
\end{defbox}

Here $n \in \N$ is identified with the set $\{0, 1, \dots, n-1\}$ (so
$\card{n} = n$ in the everyday sense), and $\card{X} = \card{n}$ means there
is a bijection $X \to n$.

\begin{defbox}
\begin{definition}[Cardinality of a finite set]
If $X$ is finite, define $\card{X} := n$, where $n \in \N$ is the unique
natural number with $\card{X} = \card{n}$.
\end{definition}
\end{defbox}

\begin{notebox}
\begin{remark}[Why $n$ is unique]
Two distinct naturals $m \neq n$ satisfy $\card{m} \neq \card{n}$ (no
bijection between sets of different sizes), so antisymmetry forces the $n$
matching $X$ to be unique. Without this fact, ``$\card{X} := n$'' would be
ill-defined.
\end{remark}
\end{notebox}

\subsection*{Counting formulas}

\begin{thmbox}
\begin{theorem}[Subtraction for subsets]
If $X$ is finite and $Y \subseteq X$, then
\[
  \card{X \setminus Y} = \card{X} - \card{Y}.
\]
\end{theorem}
\end{thmbox}

\begin{thmbox}
\begin{lemma}[Subtraction in general]
If $X$ is finite and $Y$ is finite, then
\[
  \card{X \setminus Y} = \card{X} - \card{X \cap Y}.
\]
\end{lemma}
\end{thmbox}

\begin{notebox}
\begin{remark}
The lemma reduces the general case to the previous theorem: $X \setminus Y =
X \setminus (X \cap Y)$, and $X \cap Y \subseteq X$. The point is that only
the part of $Y$ \emph{inside} $X$ can be removed; elements of $Y$ outside $X$
were never there to begin with.
\end{remark}
\end{notebox}

\subsection*{Closure properties}

\begin{thmbox}
\begin{theorem}[Intersections with finite sets]
If $X$ is finite and $Y$ is a set, then $X \cap Y$ is finite.
\end{theorem}
\end{thmbox}

\begin{cautionbox}
\begin{caution}[Board check]
The right-hand side of this theorem was partially obscured on the board.
The natural reading --- ``$X \cap Y$ is finite'' --- is the standard companion
to the next theorem, but the conclusion could instead have been a counting
formula such as $\card{X \cap Y} \leq \card{X}$. Worth confirming.
\end{caution}
\end{cautionbox}

\begin{thmbox}
\begin{theorem}[Subsets of finite sets are finite]
If $X$ is finite and $Y \subseteq X$, then $Y$ is finite.

\hfill ($\star$ exercise)
\end{theorem}
\end{thmbox}

% ============================================================
\section{Counting Formulas}
% ============================================================

With finiteness defined, we can prove the elementary counting laws on
products and unions.

\subsection*{Products}

\begin{thmbox}
\begin{theorem}[Cardinality of a product]
If $X$ and $Y$ are finite, then
\[
  \card{X \times Y} = \card{X} \cdot \card{Y}.
\]
\end{theorem}
\end{thmbox}

\begin{notebox}
\begin{remark}[Why the right-hand side makes sense]
On the natural numbers, the product is the size of the standard set
$n \cdot m = \{0, 1, 2, \dots, n \cdot m - 1\} = \{w \in \N \mid w < n \cdot m\}$,
matching the von Neumann construction used to define $\card{n} = n$. So
``$\card{X} \cdot \card{Y}$'' is genuinely a natural number, and the theorem
says the product set has exactly that many elements.
\end{remark}
\end{notebox}

\subsection*{Unions}

\begin{thmbox}
\begin{lemma}[Two-set inclusion--exclusion]
If $X$ and $Y$ are finite, then
\[
  \card{X \cup Y} = \card{X} + \card{Y} - \card{X \cap Y}.
\]
\end{lemma}
\end{thmbox}

The correction $-\card{X \cap Y}$ accounts for elements counted twice when
adding $\card{X}$ and $\card{Y}$ separately.

\begin{thmbox}
\begin{theorem}[Inclusion/Exclusion --- union bound]
If $n \in \N$ and $A_1, A_2, \dots, A_n$ are finite sets, then
\[
  \left\lvert \bigcup_{i=1}^{n} A_i \right\rvert
  \;\leq\;
  \sum_{i=1}^{n} \card{A_i}.
\]
\end{theorem}
\end{thmbox}

\begin{notebox}
\begin{remark}[Inequality, not equality]
For $n$ sets the exact formula has alternating-sign correction terms (the
full inclusion--exclusion principle); the inequality stated here drops them
and gives only the upper bound. Equality holds iff the $A_i$ are pairwise
disjoint, in which case all the correction terms vanish. The two-set lemma
above is the smallest case where the correction is actually visible:
$\card{A_1 \cup A_2} = \card{A_1} + \card{A_2} - \card{A_1 \cap A_2} \leq
\card{A_1} + \card{A_2}$.
\end{remark}
\end{notebox}

\end{document}